\section{Analysis}
\label{sec.analysis}

Throughout this section the execution is $(\epsilon,\lambda)$-typical and
\Cref{asm.majority} holds. The constants $\maxr$, $\advrone$ and $k_0$ are
those of \Cref{sec.construction}. We use one more,
\[
    \advrtwo = \frac{\advrone}{1-\epsilon} ,
\]
which is $\advrone$ measured per honest block instead of per round; it tends
to $\advr$ as $f$ and $\epsilon$ tend to $0$, and \Cref{lem.rho} shows that it
stays below $1$. We also assume $\lambda\ge\frac{2f+1}{2f(1-\advrtwo)}-2$, a
lower bound of the same order as the bound $\lambda\ge2/f$ of
\Cref{sec.model}, so that $k_0=2(\lambda+2)f\ge\frac{2f+1}{1-\advrtwo}$.

The section has four parts. \Cref{sec.growth} bounds how fast chains grow.
\Cref{sec.counting} counts the blocks on each side of a fork. \Cref{sec.safety}
proves that what the client keeps is on every honest chain, now and later.
\Cref{sec.liveness} proves that it keeps enough.

\subsection{Growth bounds}
\label{sec.growth}

The three lemmas of this subsection are routine, and their proofs are in
\Cref{appendix.analysis}. The first says that the constant $\advrtwo$ is below
$1$, so that the bounds below are well defined. The second says that honest
work on one chain arrives one block per round at most.

\begin{lemma}\label{lem.rho}
    $\advrtwo<1$.
\end{lemma}

\begin{lemma}\label{lem.honestone}
    In a typical execution, a chain contains at most one honest block computed
    in any given round.
\end{lemma}

The third lemma is the one the client relies on. It is a ceiling and a floor.
Nobody mines faster than everybody mining together, so no chain of any kind
is taller than about $\maxr fr$ at round $r$. Honest validators keep extending
the longest chain they see, so an honest chain is at least about $fr$ tall. A
chain the client is handed sits below the ceiling and has no floor, because the
adversary may deliver any chain that exists.

\begin{lemma}[Growth bounds]\label{lem.length}
    In a typical execution the following hold for every round $r\ge0$.
    \begin{enumerate}
        \item Every chain $\ch\in\bt_r$ satisfies
            $\len{\ch}-1<\maxr f\max(r,\lambda)$.
        \item If $r>\lambda$, every chain $\ch$ adopted by an honest validator
            at round $r$ satisfies $\len{\ch}-1>(1-\epsilon)f(r-1)$.
    \end{enumerate}
\end{lemma}

\subsection{Counting}
\label{sec.counting}

Take a chain $\ch$ the client holds and a chain $\ch'$ an honest validator
holds, and look at the last honest block they share. Everything $\ch$ carries
above that block and $\ch'$ does not was mined after it. In that window the
honest validators kept growing $\ch'$, one block per successful round, and
every block $\ch$ gained can be charged to a distinct adversarial success in
the same window. The lemma makes the three counts exact.

\begin{lemma}[Counting]\label{lem.count}
    In a typical execution, let $\ch\in\bt_r$ and let $\ch'$ be a chain adopted
    by an honest validator at round $r$. Let $\bst=\ch[\pos]=\ch'[\pos]$ be the
    last honest block of $\ch\meet\ch'$ and let $\rst$ be the round in which it
    was computed. Put $S=[\rst+1,r-1]$ and $S'=[\rst,r]$. Then $\rst\le r-1$,
    and
    \begin{enumerate}
        \item every block of $\ch[\pos+1{:}]$ and of $\ch'[\pos+1{:}]$ was
            computed in a round of $S'$;
        \item $\len{\ch'[\pos+1{:}]}\ge X(S)$;
        \item $\len{\ch[\pos+1{:}]}\le Z(S')+W(S)+1$.
    \end{enumerate}
\end{lemma}
\begin{proof}
    The genesis block is honest and lies on both chains, so $\bst$ exists.
    From \Cref{def.tree} and \Cref{lem.membership}, every block of $\ch$ was
    computed in a round of $[0,r]$ and every block of $\ch'$ in a round of
    $[0,r-1]$. The block $\bst$ lies on $\ch'$ and was computed at round
    $\rst$, so $\rst\le r-1$. No predictions occur, so along a chain the round
    at which a block was computed is non-decreasing in the height. Hence every
    block of $\ch[\pos+1{:}]$ and every block of $\ch'[\pos+1{:}]$ was computed
    in a round of $[\rst,r]=S'$, which is the first item.

    For the second item, $\bst$ is honest, at height $\pos$, and was computed
    at round $\rst$. From \Cref{lem.fork}, some honest validator has a chain of
    length at least $\pos+1$ at round $\rst+1$. Since $\rst+1\le r$,
    \Cref{lem.growth} from round $\rst+1$ to round $r$ gives that every honest
    validator adopts at round $r$ a chain of length at least
    $\pos+1+X([\rst+1,r-1])=\pos+1+X(S)$. In particular
    $\len{\ch'[\pos+1{:}]}=\len{\ch'}-\pos-1\ge X(S)$.

    For the third item, split the blocks of $\ch[\pos+1{:}]$ into three
    classes: $a$ adversarial blocks, $h_1$ honest blocks computed in a
    convergence opportunity of $S$, and $h_2$ remaining honest blocks. Then
    $\len{\ch[\pos+1{:}]}=a+h_1+h_2$. We bound $h_2$ first and $a+h_1$ second.

    By the first item, every honest block of $\ch[\pos+1{:}]$ was computed in a
    round of $S'$. None was computed at round $\rst$: the block $\bst$ is
    honest and was computed at round $\rst$, and by \Cref{lem.honestone} a
    chain has at most one honest block per round. So every honest block of
    $\ch[\pos+1{:}]$ was computed in a round of $S\cup\{r\}$. The $h_2$ blocks
    are the honest blocks of $\ch[\pos+1{:}]$ computed at round $r$ or in a
    successful round of $S$ that is not a convergence opportunity. From
    \Cref{lem.honestone} again, they were computed in distinct rounds. So at
    most one of them was computed at round $r$, and at most $W(S)$ of them were
    computed in successful rounds of $S$ that are not convergence
    opportunities. Hence $h_2\le W(S)+1$.

    We claim that $a+h_1\le Z(S')$. Define a map from the $a+h_1$ blocks of the
    first two classes to adversarial blocks computed in rounds of $S'$. An
    adversarial $\ch[i]$ is mapped to itself. It was computed in a round of $S'$
    by the first item. An honest $\ch[i]$ computed in a convergence
    opportunity $r'\in S$ is mapped to $\ch'[i]$. We check that $\ch'[i]$ is such
    a block. The honest validator that computed $\ch[i]$ sent a chain of length
    $i+1$ at round $r'\le r-1$, so from \Cref{lem.fork} $\len{\ch'}>i$ and
    $\ch'[i]$ exists. Since $i>\pos$ and $\ch[i]$ is honest, $\ch'[i]\ne\ch[i]$
    by the maximality of $\pos$, so $\ch'[i]$ is adversarial by
    \Cref{lem.unique}. It lies on $\ch'[\pos+1{:}]$, so by the first item
    it was computed in a round of $S'$. The map is injective: two blocks at
    distinct heights of one chain are distinct, and $\ch[i]=\ch'[i']$ would
    force $i=i'$ since no copies occur, with $\ch[i]$ adversarial and $\ch[i']$
    honest. Each adversarial block is the output of a distinct successful
    adversarial query, and $Z(S')$ counts the successful adversarial queries in
    the rounds of $S'$. Hence $a+h_1\le Z(S')$, and
    $\len{\ch[\pos+1{:}]}=a+h_1+h_2\le Z(S')+W(S)+1$.
\end{proof}

\subsection{Safety}
\label{sec.safety}

\begin{theorem}[Prefix agreement]\label{thm.prefix}
    In a typical execution, for every round $r$, every chain $\ch\in\bt_r$ and
    every chain $\ch'$ adopted by an honest validator at round $r$,
    \[
        \len{\ch\meet\ch'} > \len{\ch} - \max\left\{k_0,\
        \frac{\advrtwo\left(\len{\ch'}-\len{\ch}\right)+2f+1}{1-\advrtwo}\right\}.
    \]
\end{theorem}
\begin{proof}
    Let $\bst$, $\pos$, $\rst$, $S$ and $S'$ be as in \Cref{lem.count}, and write $d=\len{\ch[\pos+1{:}]}=\len{\ch}-\pos-1$.
    Since $\ch\meet\ch'$ contains $\bst$, we have
    $\len{\ch\meet\ch'}\ge\pos+1=\len{\ch}-d$, so it suffices to bound $d$. If
    $d<k_0$ we are done, so assume $d\ge k_0=2(\lambda+2)f$.

    We first show that $|S|>\lambda$. The subsequence $\ch[\pos+1{:}]$ consists
    of $d$ consecutive blocks of $\ch$, so from \Cref{lem.span} they were
    computed in more than $d/(2f)\ge\lambda+2$ consecutive rounds. From
    \Cref{lem.count}, all these rounds lie in $S'=[\rst,r]$, which has $|S|+2$
    elements. Hence $|S|>\lambda$, and \Cref{def.typical} applies to $S$ and to
    $S'$.

    From \Cref{lem.count,def.typical},
    \[
        d < \left(\frac{\advr}{1-f}+\epsilon\right)f(|S|+2)
        + \left(\frac{f_2}{f}+\epsilon\right)f|S| + 1 .
    \]
    From \Cref{lem.rho}, $\frac{\advr}{1-f}+\epsilon<1$, so the first term is
    smaller than $\left(\frac{\advr}{1-f}+\epsilon\right)f|S|+2f$. Therefore
    \begin{equation}\label{eq.upper}
        d < \advrone f|S| + 2f + 1 .
    \end{equation}
    From \Cref{lem.count,def.typical},
    $\len{\ch'}-\pos-1\ge X(S)>(1-\epsilon)f|S|$, so
    $f|S|<(\len{\ch'}-\pos-1)/(1-\epsilon)$ and
    \[
        d < \advrtwo(\len{\ch'}-\pos-1)+2f+1 .
    \]
    Substitute $\len{\ch'}-\pos-1=(\len{\ch'}-\len{\ch})+d$. This gives
    $d<\advrtwo(\len{\ch'}-\len{\ch})+\advrtwo d+2f+1$, and $\advrtwo<1$ by
    \Cref{lem.rho}, so $d<\frac{\advrtwo(\len{\ch'}-\len{\ch})+2f+1}{1-\advrtwo}$.
\end{proof}

The quantity $\len{\ch}-\len{\ch\meet\ch'}$ is the number of blocks of $\ch$
that are not on $\ch'$. Once it exceeds $k_0$, \Cref{thm.prefix} bounds it by
$\frac{\advrtwo}{1-\advrtwo}$ blocks for every block by which $\ch$ falls short
of $\ch'$, plus a constant, and the coefficient tends to
$\frac{\advr}{1-\advr}=\frac{t}{n-2t}$. \Cref{sec.tightness} shows that no
smaller coefficient works. The bound depends on $\len{\ch'}$, which a client
does not know. The next theorem trades that for the round, at the cost of the
gap between the ceiling and the floor of \Cref{lem.length}.

\begin{theorem}[Round form]\label{thm.round}
    In a typical execution, for every round $r$, every chain $\ch\in\bt_r$ and
    every chain $\ch'$ adopted by an honest validator at round $r$,
    \[
        \len{\ch\meet\ch'} > \len{\ch} - \max\left\{k_0,\
        \frac{\advrone\left[\maxr f(r+\lambda)-\len{\ch}+1\right]+\maxr(2f+1)}{1-f_2/f}\right\}.
    \]
\end{theorem}
\begin{proof}
    Let $\bst$, $\pos$, $\rst$ and $S$ be as in \Cref{lem.count} and write $d=\len{\ch[\pos+1{:}]}=\len{\ch}-\pos-1$. As in the proof
    of \Cref{thm.prefix}, $\len{\ch\meet\ch'}\ge\pos+1=\len{\ch}-d$, so it suffices
    to bound $d$ by the maximum in the statement, and we may assume $d\ge k_0$,
    which gives $|S|>\lambda$ and \Cref{eq.upper}.

    The prefix $\ch[{:}\pos+1]$ is a chain of $\bt_{\rst}$, because its tip
    $\bst$ was computed at round $\rst$ and no predictions occur. From
    \Cref{lem.length}, $\pos<\maxr f\max(\rst,\lambda)\le\maxr f(\rst+\lambda)$,
    hence $\rst>\frac{\pos}{\maxr f}-\lambda$. Together with
    $|S|=r-\rst-1<r-\rst$ and \Cref{eq.upper} this gives
    \[
        d < \advrone f(r+\lambda) - \frac{\advrone}{\maxr}\pos + 2f+1 .
    \]
    Put $c=\advrone/\maxr$ and observe that $c<1$: $\advrone=(1-\epsilon)\advrtwo<1$ by
    \Cref{lem.rho}, and $\maxr>1$ by its definition. Substitute $\pos=\len{\ch}-1-d$ and rearrange,
    \[
        d < \frac{\advrone f(r+\lambda)-c\,(\len{\ch}-1)+2f+1}{1-c} .
    \]
    Multiplying numerator and denominator by $\maxr$ gives
    \[
        d < \frac{\advrone\left[\maxr f(r+\lambda)-\len{\ch}+1\right]+\maxr(2f+1)}{\maxr-\advrone},
    \]
    and the definitions of $\maxr$ and $\advrone$ give
    $\maxr-\advrone=1-f_2/f$.
\end{proof}

The maximum in \Cref{thm.round} is the pruning depth $\prn(r,\len{\ch})$ of
\Cref{def.client}, so the theorem says that the client's output is a prefix of
every honest chain at the round of receipt. It does not yet say that it stays
one. Applied at a later round to the same chain, \Cref{thm.round} prunes more,
because the round grows and the chain does not. The next lemma closes the gap
with \Cref{thm.prefix} instead, which for a chain at least as long as an
honest chain prunes only a constant.

\begin{lemma}[Persistence]\label{lem.persist}
    In a typical execution, let $\ch\in\bt_r$ and let $k\ge k_0$ be an
    integer. If $\prune{\ch}{k}\preceq\ch'$ for every chain $\ch'$ adopted by
    an honest validator at round $r$, then $\prune{\ch}{k}\preceq\ch''$ for
    every chain $\ch''$ adopted by an honest validator at every round $r''>r$.
\end{lemma}
\begin{proof}
    We first note the form \Cref{thm.prefix} takes for a chain that is at
    least as long as an honest chain. Let $\ch_1$ be a chain of the block tree
    at some round and $\ch_2$ a chain adopted by an honest validator at that
    round, with $\len{\ch_1}\ge\len{\ch_2}$. The second term of the maximum in
    \Cref{thm.prefix} is then at most
    $\frac{2f+1}{1-\advrtwo}\le k_0$, so
    \begin{equation}\label{eq.longer}
        \len{\ch_1\meet\ch_2} > \len{\ch_1}-k_0 .
    \end{equation}

    Write $P=\prune{\ch}{k}$. If $k\ge\len{\ch}-1$ then $P$ has at most one
    block, so it is the genesis chain or empty and is a prefix of every chain.
    There is nothing to prove, so assume $k<\len{\ch}-1$, which gives
    $\len{P}=\len{\ch}-k$. Suppose, towards a contradiction, that the claim
    fails. Let $r'>r$ be the first round at which an honest validator $V$
    adopts a chain $\ch_B$ with $P\not\preceq\ch_B$, and let $\ch_A$ be the
    chain $V$ adopted at round $r'-1$. Then $P\preceq\ch_A$: by hypothesis if
    $r'-1=r$, by the choice of $r'$ otherwise. At round $r'$ the chain $V$ has
    is $\ch_A$ extended by at most one block. It contains $P$, so it differs
    from $\ch_B$, and since $\fn{maxvalid}$ returns a longest argument,
    $\len{\ch_B}\ge\len{\ch_A}$.

    First compare $\ch_B$ with $\ch_A$. Both $P$ and $\ch_B\meet\ch_A$ are
    prefixes of $\ch_A$ and $P\not\preceq\ch_B$, so
    $\len{\ch_B\meet\ch_A}<\len{P}$. On the other hand, $\ch_B\in\bt_{r'-1}$ by
    \Cref{lem.membership} and $\len{\ch_B}\ge\len{\ch_A}$, so \Cref{eq.longer}
    at round $r'-1$ gives $\len{\ch_B\meet\ch_A}>\len{\ch_B}-k_0\ge\len{\ch_B}-k$.
    Hence $\len{\ch_B}-k<\len{\ch}-k$, that is, $\len{\ch_B}<\len{\ch}$.

    Now compare $\ch$ with $\ch_B$. Since $\ch\in\bt_{r'}$ and
    $\len{\ch}>\len{\ch_B}$, \Cref{eq.longer} at round $r'$ gives
    $\len{\ch\meet\ch_B}>\len{\ch}-k_0\ge\len{\ch}-k=\len{P}$. Both $P$ and
    $\ch\meet\ch_B$ are prefixes of $\ch$, so
    $P\preceq\ch\meet\ch_B\preceq\ch_B$. This contradicts the choice of
    $\ch_B$.
\end{proof}

\begin{theorem}[Safety]\label{thm.safety}
    In a typical execution, if the client of \Cref{def.client} receives its
    chain at round $r$ and outputs a chain $P\ne\bot$, then $P\preceq\ch''$ for
    every chain $\ch''$ adopted by an honest validator at every round
    $r''\ge r$.
\end{theorem}
\begin{proof}
    Let $\ch$ be the chain received at round $r$ and $\hgt$ its height, so
    $P=\prune{\ch}{\prn(r,\hgt)}$ and $\ch\in\bt_r$ by \Cref{lem.membership}.
    For $r''=r$, \Cref{thm.round} gives
    $\len{\ch\meet\ch''}>\hgt-\prn(r,\hgt)\ge\len{P}$, and both $P$ and
    $\ch\meet\ch''$ are prefixes of $\ch$, so $P\preceq\ch''$. For $r''>r$, the
    integer $\lceil\prn(r,\hgt)\rceil$ is at least $k_0$ and
    $P=\prune{\ch}{\lceil\prn(r,\hgt)\rceil}$, so \Cref{lem.persist} with
    $k=\lceil\prn(r,\hgt)\rceil$ gives $P\preceq\ch''$.
\end{proof}

\begin{corollary}[Consistency]\label{cor.consistency}
    In a typical execution, let two clients, not necessarily distinct, run
    \Cref{def.client}, receive chains at rounds $r_1$ and $r_2$, and output
    $P_1\ne\bot$ and $P_2\ne\bot$. Then $P_1\preceq P_2$ or $P_2\preceq P_1$.
    In particular the output of \Cref{def.monotone} is well defined, and it is
    a prefix of every chain adopted by an honest validator at every round from
    the receipt that produced it onwards.
\end{corollary}
\begin{proof}
    Assume $r_1\le r_2$; otherwise exchange the two clients. Since $t<n$ some
    validator is honest, and it adopts a chain $\ch''$ at round $r_2$. From
    \Cref{thm.safety} applied to each client with $r''=r_2$ we get
    $P_1\preceq\ch''$ and $P_2\preceq\ch''$, and two prefixes of one chain are
    prefixes of one another. The outputs of a single client are therefore
    nested, so a longest one exists, which is the output of
    \Cref{def.monotone}. It is an output of \Cref{def.client}, so
    \Cref{thm.safety} applies to it.
\end{proof}

The two clients never talk to each other, and either may be fed every chain by
the adversary. The proof uses only that both outputs are prefixes of one honest
chain, so any two rules whose outputs satisfy \Cref{thm.safety} agree in the
same way.

\subsection{Liveness}
\label{sec.liveness}

The output is a prefix of every honest chain. This subsection says how much of
the honest chain it is, first when an honest validator sent the chain, then
when anyone did. Everything rests on the identity \Cref{eq.pref}, which we
recall: for a chain of height $\hgt$ received at round $r$,
\[
    \pref(r,\hgt) = \min\Big\{\hgt-k_0,\
    \frac{\maxr}{1-f_2/f}\big(\hgt-h_0(r)\big)\Big\},
    \qquad h_0(r)=\advrone f(r+\lambda)+2f+1 .
\]
Both branches increase with $\hgt$, so a taller chain never confirms less. The
output is not $\bot$ whenever $\pref(r,\hgt)\ge2$: then $\prn(r,\hgt)\le\hgt-2$,
so $\lceil\prn(r,\hgt)\rceil\le\hgt-1$ and the client keeps at least two
blocks.

\begin{theorem}[Honest chains]\label{thm.liveness}
    In a typical execution, let $\lambda<r'\le r$ and let the client of
    \Cref{def.client} receive at round $r$ a chain of height $\hgt$ adopted by
    an honest validator at round $r'$. Then
    \[
        \pref(r,\hgt) \ge \min\Big\{(1-\epsilon)f(r'-1)-k_0,\
        \frac{\maxr}{1-f_2/f}\big((1-\epsilon)f(r'-1)-h_0(r)\big)\Big\},
    \]
    and the output is not $\bot$ whenever the right-hand side is at least $2$.
\end{theorem}
\begin{proof}
    From \Cref{lem.length}, $\hgt>(1-\epsilon)f(r'-1)$, since $r'>\lambda$.
    Both branches of \Cref{eq.pref} increase with $\hgt$, so substituting the
    bound gives the inequality, and $\pref(r,\hgt)\ge2$ rules out $\bot$ as
    noted above.
\end{proof}

With $r'=r$ the chain is an honest chain of the current round. Since
$(1-\epsilon)-\advrone=(1-\epsilon)(1-\advrtwo)$, the second branch is then
\[
    \frac{\maxr(1-\epsilon)(1-\advrtwo)}{1-f_2/f}\,fr
\]
minus a term of order $\lambda f$, and the first branch is larger for large
$r$. So the confirmed height grows linearly in $r$, at a rate that tends to
$(1-\advr^2)f$, and a client that receives an honest chain again later confirms
strictly more. With $r'<r$ the chain is $r-r'$ rounds old. The bound falls with
the age and reaches $\bot$ when $(1-\epsilon)f(r'-1)$ drops to about $h_0(r)$,
that is, when the honest chain the client was sent is no taller than what the
adversary alone could have built by now. \Cref{cor.delay} states that
threshold. First we turn the growth rate into a confirmation time.

\begin{corollary}[Confirmation time]\label{cor.confirm}
    In a typical execution, let $\ch_1$ be a chain adopted by an honest
    validator at a round $r_1>\lambda$, and let $r>r_1$ satisfy
    $(1-\epsilon)f(r-1)\ge k_0+2$ and
    \[
        \frac{\maxr}{1-f_2/f}\big((1-\epsilon)f(r-1)-h_0(r)\big)
        \;\ge\; \maxr f(r_1+\lambda)+1 .
    \]
    Then every client that receives at round $r$ a chain adopted by an honest
    validator at round $r$ outputs a chain $P\ne\bot$ with
    $\prune{\ch_1}{k_0}\preceq P$.
\end{corollary}
\begin{proof}
    Let $\ch''$ be the chain received, of height $\hgt''$, and let $P$ be the
    output. Both branches of the bound in \Cref{thm.liveness} are at least
    $2$ under the two hypotheses, so $P\ne\bot$. From \Cref{lem.commonprefix}
    with $k=k_0\ge2\lambda f$, $\prune{\ch_1}{k_0}\preceq\ch''$, and
    $P\preceq\ch''$ by construction, so the shorter of the two is a prefix of
    the other. It remains to show
    $\len{P}\ge\len{\prune{\ch_1}{k_0}}=\len{\ch_1}-\lceil k_0\rceil$. Since
    $r>r_1$, \Cref{lem.growth} gives $\len{\ch''}\ge\len{\ch_1}$. If
    $\prn(r,\hgt'')=k_0$ then
    $\len{P}=\len{\ch''}-\lceil k_0\rceil\ge\len{\ch_1}-\lceil k_0\rceil$.
    Otherwise $\prn(r,\hgt'')>k_0$ and \Cref{eq.pref} gives
    $\pref(r,\hgt'')=\frac{\maxr}{1-f_2/f}(\hgt''-h_0(r))$. Now
    $\hgt''>(1-\epsilon)f(r-1)$ and $\len{\ch_1}<\maxr f(r_1+\lambda)+1$ by
    \Cref{lem.length}, so the second hypothesis gives $\pref(r,\hgt'')>\len{\ch_1}$,
    and $\len{P}=\hgt''+1-\lceil\prn(r,\hgt'')\rceil>\hgt''-\prn(r,\hgt'')=\pref(r,\hgt'')$. Hence
    $\len{P}>\len{\ch_1}>\len{\ch_1}-\lceil k_0\rceil$.
\end{proof}

The condition is linear in $r$ and $r_1$. It holds for every $r\ge r_1+u$
with
\[
    u = \left(\frac{1-f_2/f}{(1-\epsilon)(1-\advrtwo)}-1\right)r_1
\]
plus a term of order $\lambda$, and the coefficient tends to
$\frac{\advr}{1-\advr}=\frac{t}{n-2t}$. So a block that an honest validator
holds at round $r_1$ is confirmed by every client with an honest sender within
a $\frac{t}{n-2t}$ fraction of $r_1$ further rounds. This is the coefficient of
\Cref{thm.prefix}: a client one block behind an honest chain gives up
$\frac{t}{n-2t}$ blocks, and a block $r_1$ rounds old waits $\frac{t}{n-2t}r_1$
rounds. Both are the same exchange rate between honest and adversarial work.

\begin{corollary}[Delay]\label{cor.delay}
    In a typical execution, let $\lambda<r'\le r$ and let the client of
    \Cref{def.client} receive at round $r$ a chain $\ch$ adopted by an honest
    validator at round $r'$. If
    \[
        (1-\epsilon)f(r'-1)\ \ge\ \max\Big\{k_0+2,\
        h_0(r)+\frac{2(1-f_2/f)}{\maxr}\Big\},
    \]
    then the output is not $\bot$.
\end{corollary}
\begin{proof}
    Both branches of the bound in \Cref{thm.liveness} are then at least $2$.
\end{proof}

In the limit the condition reads $fr'\ge\advr fr$, that is $r'\ge\advr r$. An
honest chain is accepted as long as it is not older than a $(1-\advr)$
fraction of the current round, however slow the link that carried it.

\begin{corollary}[Arbitrary chains]\label{cor.floor}
    In a typical execution, let the client of \Cref{def.client} receive at
    round $r$ a chain of height $\hgt$.
    \begin{enumerate}
        \item If the output is not $\bot$, then $\hgt>k_0$ and $\hgt>h_0(r)$.
        \item $\pref(r,\hgt)\le\frac{\maxr}{1-f_2/f}\big(\hgt-h_0(r)\big)$.
    \end{enumerate}
\end{corollary}
\begin{proof}
    If the output is not $\bot$ then $\hgt-\lceil\prn(r,\hgt)\rceil\ge1$, so
    $\pref(r,\hgt)=\hgt-\prn(r,\hgt)>0$, and both branches of \Cref{eq.pref} are
    positive, which is the first item. The second item is \Cref{eq.pref}, since
    a minimum is at most its second argument.
\end{proof}

The first item is the floor. A chain no taller than $h_0(r)$, which is about
the number of blocks the adversary could have mined by itself, is rejected. The
adversary cannot make the client confirm a chain it could have built alone. The
second item is what the adversary can do instead. A chain just above $h_0(r)$
is accepted and confirms almost nothing, and that is the least an accepted
chain confirms. It is also the adversary's whole leverage over liveness. It
cannot make the client confirm anything false, by \Cref{thm.safety}, and it
cannot make the client retract anything, by \Cref{cor.consistency}, so all it
can do is hand over chains that add little. Liveness is then whatever the
honest chains that reach the client provide, which is \Cref{thm.liveness}.

\subsection{Normalised form}
\label{sec.limit}

The bounds above carry the corrections a typical execution needs. Stripping
those away shows what they say. Write $\rate=\hgt/(fr)$, as in
\Cref{sec.construction}, for the height of the chain that arrived, measured in
units of $fr$.

\begin{corollary}[Normalised form]\label{cor.limit}
    Fix $\advr$ and let a chain of height $\hgt=\rate fr$ be received at round
    $r$. Then
    \[
        \lim_{f,\epsilon\to0}\ \lim_{r\to\infty}\ \frac{\prn(r,\hgt)}{fr}
            = \advr\left(1+\advr-\rate\right),
        \qquad
        \lim_{f,\epsilon\to0}\ \lim_{r\to\infty}\ \frac{\pref(r,\hgt)}{fr}
            = (1+\advr)\left(\rate-\advr\right) .
    \]
\end{corollary}
\begin{proof}
    From \Cref{eq.pref}, $\pref(r,\hgt)/(fr)$ is the minimum of $\rate-k_0/(fr)$
    and $\frac{\maxr}{1-f_2/f}\big(\rate-h_0(r)/(fr)\big)$. As $r\to\infty$,
    $k_0/(fr)\to0$ and $h_0(r)/(fr)\to\advrone$, so the second term tends to
    $\frac{\maxr}{1-f_2/f}(\rate-\advrone)$, which is at most $\rate$ because
    $\rate<\maxr$ by \Cref{lem.length} and $\maxr-\advrone=1-f_2/f$. The
    minimum is therefore the second term. Letting $f$ and $\epsilon$ tend to
    $0$ sends $\advrone$ to $\advr$, $\maxr$ to $1+\advr$ and $f_2/f$ to $0$,
    which gives the second limit, and $\prn(r,\hgt)/(fr)=\rate-\pref(r,\hgt)/(fr)$ gives
    the first.
\end{proof}

Every number this paper reports is a value of $\rate$ in \Cref{cor.limit}, and
\Cref{tab.values} lists the four that matter. The first row is the tallest
chain that exists at round $r$, from \Cref{lem.length}, and the client keeps
all of it. The second is a chain as tall as an honest one. The third is an
honest chain of an earlier round $r'$, from \Cref{thm.liveness}. The fourth is
the floor of \Cref{cor.floor}, below which the client keeps nothing.

\begin{table}[t]
\centering
\caption{The pruning depth and the confirmed height of \Cref{cor.limit}, in
units of $fr$, at the four values of $\rate$ this paper
uses.}\label{tab.values}
\begin{tabular}{llll}
\hline
$\rate$ & the chain received & $\pref(r,\hgt)/(fr)$ & $\prn(r,\hgt)/(fr)$ \\
\hline
$1+\advr$ & the tallest chain that exists & $1+\advr$ & $0$ \\
$1$ & as tall as an honest chain & $1-\advr^2$ & $\advr^2$ \\
$r'/r$ & as tall as an honest chain of round $r'$
    & $(1+\advr)\left(\frac{r'}{r}-\advr\right)$
    & $\advr\left(1+\advr-\frac{r'}{r}\right)$ \\
$\advr$ & as tall as the adversary alone can build & $0$ & $\advr$ \\
\hline
\end{tabular}
\end{table}

The second row says two things at once. The depth $\advr^2fr$ grows with the
round, so no constant depth is sound, which is \Cref{cor.noconstant}. The
fraction $\advr^2=t^2/(n-t)^2$ does not grow at all, so the confirmed height
climbs at the steady rate $(1-\advr^2)f$. The third row prices delay. An honest
chain of round $r'$ confirms up to the height the honest chain had about
$(1+\advr)(r'-\advr r)$ rounds in, which is $(1-\advr^2)r$ for $r'=r$ and falls
by $1+\advr=\frac{n}{n-t}$ rounds for every round the chain is late. It reaches
zero at $r'=\advr r$, the fourth row, which is the delay threshold of
\Cref{cor.delay}.
